\section{Đánh giá thực
nghiệm}\label{sec:danhgia}

\subsection{Kết quả phát hiện
tamper}\label{sec:ketquaphathien}

Điều kiện sạch tạo ra verification outcome là VALID trong bảng điều khiển
xuất dữ liệu được dùng làm bằng chứng cho trạng thái không bị tamper. Các
tamper experiments tạo ra outcome non-VALID cho tất cả tamper scenarios
được định nghĩa trước. Bảng~\ref{tab:results} tóm tắt
các outcome quan sát được. Kết quả hỗ trợ H1 trong phạm vi bounded set
của các scenarios đã đánh giá.

\begin{table*}[t]
\centering
\caption{Kết quả verification quan sát được}
\label{tab:results}
\begin{tabular}{@{}p{3.6cm}p{4.4cm}p{1.3cm}@{}}
\toprule
\textbf{Scenario} & \textbf{Outcome quan sát} & \textbf{Loại} \\
\midrule
NONE & VALID & TN \\
MODIFY\_\allowbreak TRANSACTION\_\allowbreak AMOUNT & DATA\_\allowbreak TAMPERED & TP \\
DELETE\_\allowbreak TRANSACTION & RECORD\_\allowbreak COUNT\_\allowbreak MISMATCH & TP \\
INSERT\_\allowbreak FAKE\_\allowbreak TRANSACTION & RECORD\_\allowbreak COUNT\_\allowbreak MISMATCH & TP \\
MODIFY\_\allowbreak LEDGER\_\allowbreak BATCH\_\allowbreak HASH & DATA\_\allowbreak TAMPERED & TP \\
MODIFY\_\allowbreak LEDGER\_\allowbreak TRANSFORMATION & BLOCK\_\allowbreak TAMPERED & TP \\
MODIFY\_\allowbreak LEDGER\_\allowbreak PREVIOUS\_\allowbreak HASH & CHAIN\_\allowbreak BROKEN & TP \\
\bottomrule
\end{tabular}
\end{table*}
Detection rate được tính theo công thức:

\begin{equation}
Detection\ Rate = \frac{TP}{TP + FN} \times 100\%.
\end{equation}

Với sáu true positives và không quan sát thấy false negative trong các
tampered scenarios, detection rate đạt 100\% trong test set được định
nghĩa trước. Kết quả này không nên được hiểu là coverage phổ quát đối
với mọi attack vector. Nó chỉ có nghĩa là cơ chế đã phát hiện sáu tamper
vectors được lựa chọn trong thực nghiệm này.

\subsection{Kết quả traceability}\label{sec:traceability}

Verification engine trả về first broken block khi phát hiện mismatch.
Block này có thể được join với lineage events thông qua pipeline run
identifier. Lineage flow mẫu ghi nhận các transformation
Source-to-Bronze, Bronze-to-Silver và Silver-to-Gold, bao gồm input
số lượng bản ghi, output số lượng bản ghi, transformation context và status. Đối
với các tamper scenarios được đánh giá, first broken block và lineage
context khớp với affected giai đoạn ở mức giai đoạn tổng quát. Điều này hỗ trợ
H2, nhưng vẫn có giới hạn: một số giá trị giai đoạn trong source report được
ghi nhận ở mức approximate và cần được xác minh bằng SQL xuất dữ liệu trực
tiếp để có kết luận mạnh hơn.

\subsection{Kết quả chi phí phát sinh}\label{sec:overhead}

Security chi phí phát sinh được tính theo công thức:

\begin{equation}
\text{Chi phí phát sinh} = \frac{SecuredDuration - BaselineDuration}{BaselineDuration} \times 100\%.
\end{equation}

Các giá trị xấp xỉ từ bảng điều khiển cho thấy chi phí phát sinh tăng từ khoảng 160\% ở
mức 1K records lên khoảng 220\% ở mức 50K records.
Bảng~\ref{tab:overhead} tóm tắt xu hướng quan sát được.
Trường hợp 5K gần với 1K, cho thấy khả năng có ảnh hưởng của cloud
runtime variance. Nhìn chung, bằng chứng hỗ trợ H3 nhưng kèm caveats.

\begin{table*}[t]
\centering
\caption{Runtime chi phí phát sinh xấp xỉ theo số lượng bản ghi}
\label{tab:overhead}
\begin{tabular}{@{}p{2.7cm}p{5cm}p{5cm}@{}}
\toprule
\textbf{Số lượng bản ghi} & \textbf{Chi phí phát sinh xấp xỉ} & \textbf{Bằng chứng} \\
\midrule
1K & \(\sim\)160\% & Dashboard \\
5K & \(\sim\)155--160\% & Dashboard \\
10K & \(\sim\)180\% & Dashboard \\
50K & \(\sim\)220\% & Dashboard \\
\bottomrule
\end{tabular}
\end{table*}
Verification latency dao động xấp xỉ từ 24 giây đến 31 giây trong các
runs gần nhất, với một giá trị ngoại lai cao hơn. Quan sát này phù hợp với giả
định bán thực nghiệm rằng khởi động nguội (cold start) và cloud compute scheduling không
thể được kiểm soát hoàn toàn.

\subsection{Thảo luận}\label{sec:thaoluan}

Bằng chứng thực nghiệm cho thấy Hash-linked Ledger có thể làm cho hành
vi tamper sau xử lý trở nên quan sát được trong pipeline nhiều giai đoạn. Cơ
chế này đặc biệt hữu ích khi mục tiêu không phải là ngăn chặn trực tiếp
privileged writes, mà là phát hiện modification trái phép trong quá
trình verification về sau. Điều này phù hợp với thiết kế tamper-evident,
không phải tamper-proof.

Kết quả cũng làm rõ vai trò của thuật ngữ Blockchain. Nguyên mẫu sử dụng
hash chain lấy cảm hứng từ Blockchain, nhưng không triển khai
distributed consensus, replicated nút, mining, proof-of-work,
proof-of-stake hoặc smart contracts. Vì vậy, thuật ngữ chính xác hơn là
Hash-linked Tamper-evident Ledger. Sự phân biệt này quan trọng đối với
construct validity và giúp tránh overclaiming.
